\section{Results}
\label{sec:results}


\paragraph{Overview of the integrated knowledge graph}
\label{sec:knowledge_graph}

We constructed an integrated knowledge graph to analyse the representation 
and interoperability of regulatory substances across datasets. 
Figure~\ref{fig:knowledge_graph} provides an overview of the resulting 
graph structure and its main semantic layers.

\begin{figure}[!ht]
	\centering                                                             
	\includemermaid[width=1\linewidth]{images/sources/knowledge_graph.mmd}                                                                              
	\caption{
		Overview of the integrated knowledge graph constructed from 18 regulatory datasets. 
		Regulatory substances (~17,500 entities) are represented as SKOS concepts 
		and form the core layer of the graph. 
		Where possible, substances are linked to structure-based identifiers (InChIKey), 
		enabling alignment with chemical classification systems (ChemOnt) 
		and biological roles (ChEBI). 
		In addition, substances are organised into linkability tiers 
		reflecting their degree of chemical specificity, and into 
		semantic clusters derived from embedding-based similarity. 
		The graph is further enriched with annotations (~46,600 links) 
		connecting substances to regulatory sources and functional roles. 
		The figure highlights the multiple semantic layers 
		required to represent and integrate heterogeneous regulatory substance definitions.
	}\label{fig:knowledge_graph}
\end{figure} 

\paragraph{Dataset Overview}

Integration of 18 regulatory sources yielded 41,813 records. Of these,
63\% could be linked to a defined molecular structure via InChIKey.


\paragraph{Heterogeneity of Substance Definitions}
\label{sec:semantic_mismatch}

A significant share of regulatory entries (37\%) does not correspond to a defined chemical structure, consistent with Kutsarova et al.~\cite{Kutsarova2019}. These include mixtures (UVCBs), substance groups, analytical parameters, or entries without usable identifiers, reflecting differences in how “substances” are defined across regulatory contexts (Figure~\ref{fig:regulatory_groups}).

Heterogeneity varies widely between regulatory lists: some are mostly structure-based, while others rely heavily on group or mixture definitions, limiting interoperability.

Applying the seven-tier linkability taxonomy (Section~\ref{sec:linkability}) shows that 62.5\% of entries (20,909) are directly linkable via InChIKey (tier 1). Only a small fraction consists of partially resolvable groups (tiers 2–3; ${<}0.2$\%).

The remaining 37.5\% cannot be directly linked to a structure: 11.6\% have unresolved CAS numbers (tier 4), 10.8\% are mixtures or UVCBs (tier 5), ${<}0.1$\% are administrative entries (tier 6), and 14.9\% lack any usable identifier (tier 7).

Overall, more than one-third of regulatory entries cannot be mapped to a unique chemical structure, and even CAS numbers do not guarantee unambiguous identification. Consequently, structure-based integration applies only to part of the data, requiring alternative approaches for the remainder.

\begin{figure}[H]
	\centering
	\includegraphics[width=\linewidth]{Analysis_4d_Linkability_taxonomy_general_overview.pdf}
	\caption{
		Seven-tier linkability taxonomy applied to all 33,452 unique regulatory substance entries. 
		Tier 1 (62.5\%) corresponds to substances that can be directly linked to a defined chemical structure via InChIKey. 
		The remaining tiers represent decreasing levels of structural specificity, including group definitions, mixtures, and entries lacking actionable identifiers. 
		Overall, 37.5\% of regulatory substances cannot be resolved to a unique chemical structure.
	}
	\label{fig:analysis4d}
\end{figure}


\paragraph{Limitations of CAS Numbers}~\ref{sec:question1}

Although CAS numbers are widely used, they do not provide a reliable
unique identifier. Multiple CAS numbers may refer to the same structure,
and conversely, a single CAS number may correspond to multiple structures.
This confirms that CAS-based linking is insufficient for interoperability.

\paragraph{Limited Interoperability Across Regulatory Lists}~\ref{sec:question2}

Cross-list analysis shows that most substances are unique to a single
regulatory source. Overlap between lists is limited, reflecting the
domain-specific design of regulatory instruments rather than a shared
chemical reference framework.

\paragraph{Failure of Semantic Approaches for Non-Structure Substances}~\ref{sec:question3}
Embedding-based clustering of non-structure substance names yields semantically coherent groups, but does not resolve chemical identity. While clusters reflect lexical or contextual similarity between substance descriptions, they do not correspond to well-defined molecular entities. 

\begin{figure}[H]
	\centering
	\includegraphics[width=\linewidth]{Analysis_9n_quality_landscape.pdf}
	\caption{Match quality landscape: score gap $\Delta_{12}$ versus best match score $s_1$ for all non-structure substance--ChemOnt class label pairs. Entries in the upper-right region satisfy both quality thresholds and are retained as high-confidence matches.}
	\label{fig:analysis9n}
\end{figure}




To evaluate whether semantic similarity can recover chemical meaning, we attempted to map non-structure substances to ChemOnt taxonomy classes. Fewer than 1\% yielded high-confidence matches. Among the top 100 ranked results, the average confidence was 0.55, and manual inspection showed that only five exactly matched ChemOnt classes. This indicates that semantic approaches are largely insufficient for reconstructing chemical identity.

The match quality landscape (Figure~\ref{fig:analysis9n}) plots the score gap $\Delta_{12}$ between the best and second-best matches against the best match score $s_1$ for all substance--class pairs. Entries in the upper-right region (high $s_1$, high $\Delta_{12}$) correspond to unambiguous matches and represent the 107 high-confidence mappings retained in our analysis. In contrast, entries with high $s_1$ but low $\Delta_{12}$ indicate genuine ambiguity between multiple candidate classes, reflecting the underspecified nature of the original substance definitions.

The predominance of low-confidence and ambiguous matches demonstrates that semantic and embedding-based methods cannot compensate for the absence of structure-based identifiers. In particular, these approaches fail systematically for substance groups, mixtures, and analytical parameters, where no single underlying chemical structure exists.





\paragraph{Scale of the Interoperability Problem}
For the more than 15,000 non-structure entries, full manual reconciliation would require on the order of hundreds of millions of pairwise comparisons. Under conservative assumptions, this corresponds to several thousand person-years of expert effort.

\label{sec:workload}
\begin{figure}[!ht]
	\centering
	\includegraphics[width=\linewidth]{Analysis_10c_workload_non_structured.pdf}
	\caption{
		Estimated pairwise curation workload for 11,250 non-structure entries across the combined European regulatory dataset, broken down by source. 
		The total corresponds to approximately 630,000 person-days (3,150 person-years) under an assumed throughput of 100 pair assessments per person-day. 
		The quadratic growth of the workload illustrates that manual reconciliation is not scalable, reinforcing the need for structure-based and ontology-driven representations at the source.
	}
	\label{fig:analysis10c}
\end{figure}



Figure~\ref{fig:analysis10c} quantifies this workload across all regulatory sources, illustrating both the absolute scale and the quadratic growth that makes each additional non-structure entry disproportionately costly. Even moderate increases in the number of unresolved entries lead to a rapidly escalating curation burden.

Manual pairwise reconciliation at this scale is not operationally feasible. This result highlights a fundamental constraint: interoperability cannot be achieved through retrospective harmonisation of heterogeneous substance definitions.

Instead, the only scalable strategy is to address the problem at its source. By requiring structure-resolvable identifiers and ontology-based representations in regulatory definitions wherever chemically applicable, the number of entries requiring manual reconciliation can be structurally reduced rather than managed post hoc.


\paragraph{Knowledge Graph Integration}

Structure-defined substances can be successfully integrated across
regulatory lists, chemical taxonomies (ChemOnt), and biological
annotations (ChEBI) using InChIKey as a linking key. This demonstrates
the feasibility of a linked-data approach when structure information is
available.



\paragraph{Priority Patterns}

The integrated knowledge graph enables the combination of regulatory 
and biological annotations into a shared analytical space. 
Using the composite priority score defined in Section~\ref{sec:priority_score}, 
we identify substances and chemical classes that are prominent across 
both regulatory presence and biological role annotations.

The bubble chart (Figure~\ref{fig:analysis14e}) visualises this 
priority space by plotting $\text{list\_score}$ against 
$\text{bio\_score}$ for all scored substances. 
Bubble size is proportional to the number of regulatory source memberships. 


\begin{figure}[!ht]
	\centering
	\includegraphics[width=\linewidth]{Analysis_14e_Bubble_priority.pdf}
	\caption{Bubble chart of composite priority space: regulatory list score versus biological hazard score for all scored substances. Bubble size is proportional to the number of regulatory source memberships. Substances in the upper-right quadrant combine broad regulatory coverage with documented biological hazard evidence.}
	\label{fig:analysis14e}
\end{figure}

Substances located in the upper-right quadrant combine high regulatory 
coverage with multiple biological role annotations, reflecting their 
prominence across heterogeneous data sources. 
Examples include bisphenol~A (BPA) and dibutyl phthalate, which appear 
consistently across regulatory lists and are associated with multiple 
biological roles in ChEBI.

Importantly, this analysis does not constitute a toxicological ranking, 
but demonstrates how interoperable data representations enable the 
integration of regulatory and ontology-based information within a 
single analytical framework.

At the class level, a small number of chemical classes accounts for 
a large fraction of regulated substances. This concentration suggests 
that class-based representations can provide a scalable bridge between 
individual substances and regulatory group definitions, supporting 
alignment across regulatory contexts.




